\chapter{IMC Prosperity: Competition, Mechanics, Tooling}
\label{ch:prosperity_setup}
\chaptermark{Prosperity: competition and tooling}

You are about to compete in IMC Prosperity.  Everything covered so far
applies directly: limit
order books (\cref{ch:lob}), market making (\cref{ch:market_making}),
statistical arbitrage (\cref{ch:basket_etf_arb}), options pricing
(\cref{ch:black_scholes,ch:greeks_hedging,ch:vol_surface}), and the
``which strategy when'' framework (\cref{ch:which_strategy_when}).  This chapter is the operational briefing: competition
structure, Python API, tooling, and the product archetypes that recur
every year.

\begin{backgroundbox}[Prerequisites]
\begin{itemize}
  \item \Cref{ch:lob}: limit order books, price--time priority,
  bid/ask/spread, order types.
  \item \Cref{ch:market_making}: the Avellaneda--Stoikov framework,
  inventory risk, quoting around fair value.
  \item \Cref{ch:which_strategy_when}: the four-layer decision
  pipeline mapping market observables to strategies.
\end{itemize}
\end{backgroundbox}

%======================================================================
\section{What is IMC Prosperity?}
\label{sec:prosperity_what}
%======================================================================

IMC Prosperity is a global algorithmic trading competition run by IMC
Trading, a major market-making firm, for university
students~\citep{prosperity4wiki}.  Teams of up to four submit Python
algorithms that trade a simulated exchange populated by bots.  The
objective: \emph{maximise PnL} measured in the in-game currency
(SeaShells in P3, XIRECs in P4).

The competition spans $\approx$16 days in five rounds.  Each round
introduces 1--3 new products; earlier-round products persist.  Alongside
each algorithmic round is a \textbf{manual challenge}, a human-decision
task (optimisation under uncertainty, game theory) scored independently.
Combined PnL across all rounds sets the leaderboard.

Scale: P3 (2025) attracted over 12{,}000 teams; P4 is expected to be
similar.  Ranking is pure PnL; there are no style points.

\begin{keyfactbox}[Prosperity at a glance]
\begin{itemize}
  \item \textbf{Organiser:} IMC Trading.
  \item \textbf{Format:} 5 algorithmic rounds + 5 manual rounds over
  $\approx$16 days.
  \item \textbf{Language:} Python (submitted as a single
  \texttt{.py} file).
  \item \textbf{Execution:} your algorithm trades independently
  against bots; there is \emph{no} direct interaction between
  different teams' algorithms.
  \item \textbf{Scoring:} total PnL (realised + unrealised) across all
  rounds.  Algo and manual PnL are summed for the overall leaderboard.
  \item \textbf{Scale:} 10{,}000--12{,}000+ teams.
\end{itemize}
\end{keyfactbox}

Prosperity compresses the core challenges of systematic trading (fair
value, inventory, spread capture, stat arb, options) into a
fast-feedback environment where iteration takes hours, not months.
Every strategy archetype in Parts~III--V has appeared in at least one
round.

%======================================================================
\section{Competition structure}
\label{sec:prosperity_structure}
%======================================================================

\subsection{Five rounds, cumulative products}

Rounds~1--2 last 72\,h each; Rounds~3--5 last 48\,h each.  A four-day
intermission separates Round~2 from~3.  A $\approx$3-hour scoring window
between rounds processes Round $n$'s PnL while Round $n{+}1$ is
locked~\citep{prosperity4wiki}.  \Cref{tab:round_schedule} shows the P4
schedule; the scoring window is dead time for submission but valuable
for analysis.

\begin{table}[htbp]
\centering
\small
\caption{Prosperity~4 round schedule (all times CEST / UTC+2).}
\label{tab:round_schedule}
\begin{tabular}{@{}llll@{}}
\toprule
\textbf{Round} & \textbf{Start} & \textbf{Close} &
\textbf{Duration} \\
\midrule
Tutorial & (signup) & 14 Apr 12:00 & variable \\
Round 1 & 14 Apr 12:00 & 17 Apr 12:00 & 72\,h \\
Round 2 & 17 Apr 12:00 & 20 Apr 12:00 & 72\,h \\
\emph{Intermission} & 20 Apr 12:00 & 24 Apr 12:00 & 4 days \\
Round 3 & 24 Apr 12:00 & 26 Apr 12:00 & 48\,h \\
Round 4 & 26 Apr 12:00 & 28 Apr 12:00 & 48\,h \\
Round 5 & 28 Apr 12:00 & 30 Apr 12:00 & 48\,h \\
\bottomrule
\end{tabular}
\end{table}

Each round boundary adds 1--3 new products; all earlier products remain
tradable.  By P3 Round~5, the universe had grown from 3 to 15 products.
This rewards modular code: a well-factored Round~1 strategy should keep
running unmodified while new ones are bolted on.

\begin{keyfactbox}[Prosperity round structure]
\begin{enumerate}
  \item Each round introduces 1--3 new products; prior products persist.
  \item Algo and manual submissions are locked at round close; the
  last \emph{active} submission is used.
  \item PnL = realised + unrealised at the end of the round's
  simulation.
  \item The simulation runs for 1{,}000 iterations during testing
  and 10{,}000 iterations for the official scoring run.
  \item Position limits are enforced \emph{per product} by the
  exchange.  Exceeding the limit causes \emph{all} orders for that
  product to be cancelled in that iteration.
\end{enumerate}
\end{keyfactbox}

\subsection{Algorithmic vs.\ manual rounds}

Each round has two independent components:

\begin{enumerate}
  \item \textbf{Algorithmic challenge.}  Upload a Python file containing
  your \texttt{Trader} class.  Unlimited revisions; only the last
  \emph{active} submission at round close is scored.  Historical price
  and trade CSVs for each new product let you study dynamics before
  coding.

  \item \textbf{Manual challenge.}  A human-decision task: optimisation
  under uncertainty, game theory, news trading.  Manual PnL is
  \emph{added} to algo PnL.  Prior tasks have ranged from triangular FX
  arbitrage to Nash-equilibrium container
  selection~\citep{frankfurthedgehogs2025}.
\end{enumerate}

The manual round is inactive during the tutorial period.  Skipping the
manual round leaves free money on the table; even a mediocre submission
beats zero.

\subsection{PnL and scoring}

Round PnL = realised (closed trades) + unrealised (mark-to-market on
open positions) at simulation end.  Concretely, realised PnL sums
$-\sum p_i q_i$ over your fills (signed: $q_i>0$ for buys costs cash,
$q_i<0$ for sells earns cash), and unrealised PnL marks any leftover
position $q$ at the final mid-price $m$ as $q \cdot m$, so a long
position counts at what it could be sold for, a short at what it would
cost to cover.  Currency: XIRECs (P4) or SeaShells (P3).  The overall
leaderboard ranks by \emph{total} PnL across rounds (algo + manual),
with sub-leaderboards for each.

A subtlety: since PnL is cumulative and products persist, a later round
losing money on a Round~1 product reduces your \emph{total} score.  If
you cannot improve a product's strategy, sometimes disabling it beats
letting it bleed.

\subsection{Position limits}

Each product has a position limit $L$: your signed position $q$ must
satisfy $-L \le q \le L$.  Enforcement is aggressive: if the aggregate
of your buy (sell) quantities plus $q$ would exceed $+L$ ($-L$),
\emph{every} order for that product in that iteration is rejected, not
just the excess.

Typical P3 limits ranged from 50 (simple MM products like Rainforest
Resin, Kelp) to 400 (Volcanic Rock, the options underlying).  Limits
are disclosed in the product spec at round start.

\begin{pitfallbox}[Position limit enforcement is all-or-nothing]
There is no partial rejection: \emph{all} orders for the product are
voided if the aggregate could push past the limit.  Compute remaining
capacity up front:
\[
  \text{max\_buy\_qty} = L - q, \qquad
  \text{max\_sell\_qty} = L + q.
\]
Safe pattern: compute these bounds at the start of \texttt{run()} and
thread them through every order-building function.
\end{pitfallbox}

\subsection{Year-over-year product archetypes}
\label{sec:archetypes}

Names and narratives change each year, but the product archetypes
recur.  Recognising the archetype in the first hour lets you skip ``what
is this?'' and go straight to the matching strategy.
\Cref{tab:archetypes} maps archetypes across P1--P4.

\begin{table}[htbp]
\centering
\small
\caption{Recurring product archetypes across Prosperity editions.  The
``Strategy chapter'' column points to where the matching technique is
derived in this book.}
\label{tab:archetypes}
\begin{tabular}{@{}p{2.8cm}p{2.0cm}p{2.0cm}p{2.4cm}p{1.8cm}p{2.4cm}@{}}
\toprule
\textbf{Archetype} & \textbf{P1} & \textbf{P2} & \textbf{P3} &
\textbf{P4} & \textbf{Strategy ch.} \\
\midrule
Flat-value market making
  & Amethysts
  & Amethysts
  & Rainforest Resin
  & Emeralds\textsuperscript{*}
  & \cref{ch:market_making} \\[4pt]
Trending / random-walk MM
  & ---
  & Starfruit
  & Kelp
  & Tomatoes\textsuperscript{*}
  & \cref{ch:market_making,ch:mean_reversion_ou} \\[4pt]
Signal-driven (informed trader)
  & ---
  & ---
  & Squid Ink
  & TBD
  & \cref{ch:momentum,ch:ml_trading} \\[4pt]
Basket / ETF stat arb
  & ---
  & Gift Baskets
  & Picnic Baskets (2)
  & TBD
  & \cref{ch:basket_etf_arb} \\[4pt]
Options / vol surface
  & ---
  & ---
  & Volcanic Rock Vouchers
  & TBD
  & \cref{ch:black_scholes,ch:vol_surface} \\[4pt]
Cross-exchange / locational
  & ---
  & Orchids
  & Macarons
  & TBD
  & \cref{ch:arbitrage_zoo} \\
\bottomrule
\multicolumn{6}{@{}l}{\footnotesize \textsuperscript{*}Tutorial-round
products; Round~1 onward may introduce additional products.}
\end{tabular}
\end{table}

Every year has at least one flat-value MM product, one trending MM
product, one multi-leg stat-arb product, and (since P3) one options
product.  Cross-exchange arbitrage with conversion fees has appeared
every edition since P2.

\subsubsection*{How to identify the archetype on day 1}

Run this checklist when a new product appears:

\begin{enumerate}
  \item \textbf{Plot prices.}  Flat (constant fair), trending (random
  walk), or mean-reverting?
  \item \textbf{Compute spread.}  Wide ($>$4 ticks) $\Rightarrow$ MM
  opportunity; tight ($\le$2) $\Rightarrow$ edge elsewhere.
  \item \textbf{Related products.}  A basket alongside constituents
  $\Rightarrow$ basket stat-arb: compute synthetic and plot residual.
  \item \textbf{Conversion observations.}  \texttt{observations} with
  bid/ask/tariff for an external market $\Rightarrow$ cross-exchange.
  \item \textbf{Option-like.}  Multiple strike variants alongside an
  underlying $\Rightarrow$ options; compute IV for each.
  \item \textbf{Trade flow.}  \texttt{market\_trades} filtered by size
  and direction revealing predictable bots (e.g.\ always buying the
  daily low) $\Rightarrow$ signal-driven.
\end{enumerate}

\begin{intuitionbox}
Each new Prosperity product is an instance of $\approx$6 archetypes.
Day-1 step: \emph{plot price, compute spread, ask ``which archetype?''}
Then pull the matching skeleton from the relevant chapter and adapt
parameters.  With notebooks prepared in advance, the checklist should
take under 15 minutes per product.
\end{intuitionbox}

%======================================================================
\section{The Python API}
\label{sec:prosperity_api}
%======================================================================

Your algorithm is a single Python file containing a \texttt{Trader}
class.  The simulation calls \texttt{Trader.run()} once per iteration
(``tick''), passing the current state and expecting back an order dict.
This section documents every field you will touch.

\subsection{The \texttt{Trader} class}

The minimal skeleton is:

\begin{lstlisting}[caption={Minimal \texttt{Trader} skeleton.},
  label={lst:trader_skeleton}]
from datamodel import OrderDepth, TradingState, Order

class Trader:
    def run(self, state: TradingState) -> tuple:
        result = {}
        for product in state.order_depths:
            orders = []
            # --- strategy logic here ---
            result[product] = orders
        traderData = ""   # JSON string for persistence
        conversions = 0   # conversion requests (if any)
        return result, conversions, traderData
\end{lstlisting}

\texttt{run} takes a \texttt{TradingState} and returns a triple
\texttt{(result, conversions, traderData)}:

\begin{itemize}
  \item \texttt{result}: \texttt{dict[str, list[Order]]} mapping product
  \texttt{Symbol} to orders.
  \item \texttt{conversions}: integer conversion request (cross-exchange
  products only; else~0).
  \item \texttt{traderData}: string delivered back as
  \texttt{state.traderData} next call, your only persistence mechanism
  (see pitfall below).
\end{itemize}

Return within 900\,ms (typical: $<$100\,ms); a timeout produces no
orders that iteration.

\subsection{The \texttt{TradingState} class}
\label{sec:tradingstate}

\texttt{TradingState} is the full world snapshot visible to your
algorithm:

\begin{table}[htbp]
\centering
\small
\caption{\texttt{TradingState} fields.}
\label{tab:tradingstate}
\begin{tabular}{@{}lll@{}}
\toprule
\textbf{Field} & \textbf{Type} & \textbf{Meaning} \\
\midrule
\texttt{timestamp} & \texttt{int} &
  Simulation clock (increases each iteration). \\
\texttt{traderData} & \texttt{str} &
  The string you returned last iteration. \\
\texttt{listings} & \texttt{dict[str, Listing]} &
  Product metadata (symbol, denomination). \\
\texttt{order\_depths} & \texttt{dict[str, OrderDepth]} &
  Current order book per product. \\
\texttt{own\_trades} & \texttt{dict[str, list[Trade]]} &
  Trades \emph{your} algorithm did since last tick. \\
\texttt{market\_trades} & \texttt{dict[str, list[Trade]]} &
  Trades \emph{other} bots did since last tick. \\
\texttt{position} & \texttt{dict[str, int]} &
  Your signed position per product. \\
\texttt{observations} & \texttt{Observation} &
  External data (sunlight, tariffs, etc.). \\
\bottomrule
\end{tabular}
\end{table}

The two heavy-use fields are \texttt{order\_depths} (what you can trade
against) and \texttt{position} (inventory).  \texttt{own\_trades}
confirms last-iteration fills; \texttt{market\_trades} reveals bot
behaviour and sometimes identifies specific bot personalities
(\cref{ch:algo_rounds}).

\subsection{The \texttt{OrderDepth} class}
\label{sec:orderdepth}

\texttt{OrderDepth} \emph{is} the LOB snapshot from \cref{ch:lob}, as
two dictionaries:

\begin{lstlisting}[caption={\texttt{OrderDepth} class definition.},
  label={lst:orderdepth}]
class OrderDepth:
    def __init__(self):
        self.buy_orders: dict[int, int] = {}
        self.sell_orders: dict[int, int] = {}
\end{lstlisting}

Both map integer \textbf{price} to integer \textbf{total volume} at
that level: \texttt{buy\_orders} holds resting bid levels,
\texttt{sell\_orders} holds resting ask levels.

\begin{pitfallbox}[\texttt{OrderDepth} sign convention]
\texttt{sell\_orders} values are \textbf{negative}.  4 lots offered at
price~11 is \texttt{\{11: -4\}}, not \texttt{\{11: 4\}}.  Negate to get
a positive buy quantity:
\begin{lstlisting}[numbers=none, frame=none, backgroundcolor=\color{red!3}]
best_ask, best_ask_vol = list(od.sell_orders.items())[0]
buy_qty = -best_ask_vol  # now positive
\end{lstlisting}
Nearly every team hits this once.  If your first submission produces
zero trades, check this first.
\end{pitfallbox}

Prices are integers and the book is uncrossed (every bid $<$ every
ask), so:

\begin{lstlisting}[caption={Extracting best bid/ask and mid price.},
  label={lst:bestbidask}]
best_bid = max(od.buy_orders.keys())
best_ask = min(od.sell_orders.keys())
mid = (best_bid + best_ask) / 2
\end{lstlisting}

The discrete-price analogue of $\spreads = \askp - \bidp$
(\cref{ch:lob}).

\subsection{The \texttt{Order} class}

Construct an \texttt{Order} to submit:

\begin{lstlisting}[caption={\texttt{Order} class usage.},
  label={lst:order_class}]
class Order:
    def __init__(self, symbol: str, price: int,
                 quantity: int) -> None:
        self.symbol = symbol
        self.price = price
        self.quantity = quantity
\end{lstlisting}

\begin{itemize}
  \item \texttt{quantity > 0} $\Rightarrow$ \textbf{buy} order.
  \item \texttt{quantity < 0} $\Rightarrow$ \textbf{sell} order.
\end{itemize}

If your buy price meets or exceeds a resting sell, the trade executes
immediately at the \emph{resting} price (price--time priority, as in
\cref{ch:lob}).  Unmatched residuals rest in the book; bots may trade
against them before the next tick, and if none do, they are cancelled
at iteration end.

\subsection{The \texttt{Trade} class}

Each \texttt{own\_trades}/\texttt{market\_trades} element is a
\texttt{Trade} with \texttt{symbol, price, quantity, buyer, seller,
timestamp}.  Counterparty names are hidden: \texttt{buyer}/\texttt{seller}
are empty strings unless \emph{your} algorithm is one of the parties (in
which case \texttt{"SUBMISSION"}).  In P3 Round~5, trader IDs were
briefly revealed, enabling direct bot identification.

\subsection{Matching and execution}

Matching uses price--time priority (\cref{ch:lob}):
\begin{enumerate}
  \item Zero latency: you never lose a race to a bot.
  \item Buys match cheapest sells first; sells match most expensive buys
  first.
  \item Fills at the \emph{resting} price, so you may get a better price
  than quoted.
  \item Unmatched residuals rest; if no bot trades, they are cancelled
  before the next iteration.
\end{enumerate}

This ``snapshot-then-act'' model has no latency, no queue priority, and
full book visibility before acting.  Strategic consequence:
\emph{information extraction} dominates \emph{speed}.

\subsection{Conversions and observations}

Some products support \textbf{conversions}: trading with a foreign
exchange at known prices plus fees (the cross-exchange archetype,
\cref{sec:archetypes}).  When conversions are available,
\texttt{observations} contains a \texttt{ConversionObservation} with
foreign bid/ask, transport fees, and import/export tariffs.  Execute by
returning a non-zero integer as the \texttt{conversions} element.

A per-tick conversion limit (e.g.\ 10 units for P3 Macarons) caps
throughput.  You must hold sufficient inventory; the exchange will not
create positions from nothing.

\subsection{Worked example: two iterations of trading}
\label{sec:worked_iterations}

Consider \texttt{PRODUCT1} with position limit~10, estimated fair value
13, and current position~3.

\textbf{Iteration 1.}  The \texttt{TradingState} arrives with:
\begin{lstlisting}[numbers=none, frame=none,
  backgroundcolor=\color{codebg}]
order_depths["PRODUCT1"] = OrderDepth(
    buy_orders  = {10: 7, 9: 5},
    sell_orders = {11: -4, 12: -8}
)
position["PRODUCT1"] = 3
\end{lstlisting}

Both sells (11, 12) are below fair value~13, so buy.  Max buy capacity
is $L - q = 10 - 3 = 7$ lots.  Send:
\begin{lstlisting}[numbers=none, frame=none,
  backgroundcolor=\color{codebg}]
Order("PRODUCT1", 12, 7)   # buy up to 7 at price 12
\end{lstlisting}

The exchange matches: 4 lots at~11 (you pay 11, not 12, since fills
occur at the resting price), then 3 of 8 lots at~12.  You bought 7 lots at average
$(4 \times 11 + 3 \times 12)/7 \approx 11.43$.

\textbf{Iteration 2.}  The next \texttt{TradingState} shows:
\begin{lstlisting}[numbers=none, frame=none,
  backgroundcolor=\color{codebg}]
own_trades["PRODUCT1"] = [
    Trade("PRODUCT1", 11, 4, buyer="SUBMISSION"),
    Trade("PRODUCT1", 12, 3, buyer="SUBMISSION"),
]
position["PRODUCT1"] = 10   # at the limit
\end{lstlisting}

At the limit, any buy causes \emph{all} \texttt{PRODUCT1} orders to be
rejected.  Options: sell or do nothing.

\begin{intuitionbox}
Prosperity's iteration model is a discrete-time LOB (\cref{ch:lob}):
each tick gives a fresh snapshot, you submit, you observe fills.  No
race condition, no latency advantage, no queue priority.  The edge is
\emph{fair-value estimation} and \emph{inventory management} against
position limits.
\end{intuitionbox}

%======================================================================
\section{Tooling}
\label{sec:prosperity_tooling}
%======================================================================

Top teams backtest locally, visualise order books, and iterate.  This
section covers the de~facto standard tools.

\subsection{\texttt{prosperity3bt}: the community backtester}

Jasper Merle's open-source backtester,
\texttt{prosperity3bt}~\citep{jmerle_backtester}, replays historical
data against your \texttt{Trader} locally, producing logs in the
official submission format.

\begin{lstlisting}[language=bash, caption={Installing and running the
  backtester.}, label={lst:backtester}]
# Install
pip install -U prosperity3bt

# Run on all days from round 1
prosperity3bt trader.py 1

# Run on a specific day (round 1, day 0)
prosperity3bt trader.py 1-0

# Merge PnL across days
prosperity3bt trader.py 1 --merge-pnl

# Open result in the visualiser automatically
prosperity3bt trader.py 1 --vis
\end{lstlisting}

The backtester matches against recorded order depths and market trades,
enforcing position limits identically to the official environment.

\begin{pitfallbox}[Backtester vs.\ official environment]
Replayed bots do not \emph{react} to your orders.  Therefore:
\begin{itemize}
  \item Strategies taking existing liquidity (buying underpriced asks):
  backtester is accurate.
  \item Strategies relying on bots trading against \emph{your} resting
  quotes (passive MM): backtester over- or understates fills depending
  on \texttt{-{}-match-trades}.
\end{itemize}
Use the official test environment for strategies where bot interaction
matters.
\end{pitfallbox}

\subsection{Jasper Merle's visualiser}

The P3 Visualiser~\citep{jmerle_visualiser} parses backtester logs and
displays:
\begin{itemize}
  \item Order-book levels over time (bids/asks with depth).
  \item Your trades overlaid on price.
  \item PnL and position curves per product.
\end{itemize}

\texttt{prosperity3bt trader.py 1 -{}-vis} opens it after a backtest.
The Frankfurt Hedgehogs built a custom dashboard (indicator
normalisation, trader-ID filtering), but the community visualiser is
sufficient for most teams.

\subsection{Recommended repository layout}

A clean structure prevents confusion as rounds add products:

\begin{lstlisting}[language=bash, numbers=none, frame=single,
  caption={Suggested repository layout.}, label={lst:repo_layout}]
prosperity/
  src/
    trader.py         # Your Trader class (submitted)
    datamodel.py      # Copy of the official datamodel
    strategies/       # One module per product archetype
      market_making.py
      basket_arb.py
      options.py
  logs/               # Backtester output logs
  analysis/           # Jupyter notebooks, plots
  data/               # Sample CSV data from the dashboard
  README.md
\end{lstlisting}

Keep \texttt{trader.py} a thin dispatcher that imports strategy
modules, so new products add without touching existing code.  The
Frankfurt Hedgehogs used a \texttt{ProductTrader} base class with
per-product subclasses sharing utilities for book parsing, position
tracking, and logging~\citep{frankfurthedgehogs2025}.

\subsection{Debugging workflow}

The platform returns a log file with all \texttt{print()} output from
your \texttt{run()}.  Workflow:

\begin{enumerate}
  \item \textbf{Structured logging.}  Print timestamp, product,
  position, best bid/ask, and orders sent.  Compact format for scanning
  thousands of iterations.
  \item \textbf{Run locally first.}  \texttt{prosperity3bt trader.py 1
  -{}-print} streams output in real time.
  \item \textbf{Visualise.}  Feed the log into the visualiser.  Overlay
  trades on the book.  Look for obvious mistakes: wrong side, quotes far
  from fair, flat PnL when you expect activity.
  \item \textbf{Validate on the platform.}  Upload to the official test
  environment (1{,}000-iter run).  Large divergence from local PnL
  usually indicates bot-interaction effects.
\end{enumerate}

\subsection{Historical data}

For each new product the platform provides several days of CSVs:
\begin{itemize}
  \item \textbf{Prices CSV:} order-book state (all levels with volumes)
  at every timestamp.
  \item \textbf{Trades CSV:} every trade, with price, quantity, and (in
  later rounds) buyer/seller IDs.
\end{itemize}

The backtester ships with data for all rounds.  Use these CSVs for
exploratory analysis (plot prices, compute spreads, check mean
reversion, spot bot patterns) \emph{before} writing strategy code.

\subsection{AWS Lambda constraints}

The official environment runs on AWS Lambda.  Two constraints:

\begin{enumerate}
  \item \textbf{Memory:} $\approx$256\,MB.  \texttt{numpy}, \texttt{pandas},
  \texttt{statistics} are available, but heavy structures hit the ceiling.
  \item \textbf{Statelessness:} globals and class attributes are not
  guaranteed to persist across \texttt{run()} calls.
\end{enumerate}

\begin{pitfallbox}[No persistent state between ticks]
Lambda can recycle the container between iterations, so instance
variables (\texttt{self.history}), class variables, and module globals
are \textbf{not} reliable across \texttt{run()} calls.  The only
reliable persistence is the \texttt{traderData} string returned as the
third tuple element; it is delivered back as \texttt{state.traderData}
next call.

Serialise state with \texttt{json} (or \texttt{jsonpickle}):
\begin{lstlisting}[numbers=none, frame=none,
  backgroundcolor=\color{red!3}]
import json

class Trader:
    def run(self, state: TradingState) -> tuple:
        # Restore state
        data = json.loads(state.traderData) \
            if state.traderData else {}
        history = data.get("prices", [])

        # ... strategy logic using history ...
        mid = (best_bid + best_ask) / 2
        history.append(mid)

        # Save state (max 50,000 characters)
        traderData = json.dumps({"prices": history})
        return result, conversions, traderData
\end{lstlisting}
\texttt{traderData} is truncated to 50{,}000 characters, so keep state
compact.
\end{pitfallbox}

\subsection{Supported libraries}

Python~3.12, with:
\begin{itemize}
  \item \texttt{numpy}, \texttt{pandas}: numerics and data.
  \item \texttt{statistics}, \texttt{math}: stdlib numerics.
  \item \texttt{jsonpickle}: serialisation for \texttt{traderData}.
  \item All Python~3.12 stdlib (\texttt{collections}, \texttt{itertools},
  \texttt{functools}, \dots).
\end{itemize}
No other libraries.  Notably \texttt{scipy} is \emph{not} available, so
for a normal CDF (Black--Scholes) use \texttt{statistics.NormalDist} or
implement it yourself.

%======================================================================
\section{The simulation model and bot ecosystem}
\label{sec:prosperity_sim}
%======================================================================

Understanding the simulation internals is essential for strategies that
translate from backtest to live.

\subsection{Iteration flow}

Each iteration proceeds in a fixed
sequence~\citep{frankfurthedgehogs2025}:
\begin{enumerate}
  \item All previous player orders cancelled.
  \item Deep-liquidity MM bots post new quotes (the ``walls'').
  \item Taker bots may cross and generate trades.
  \item Your \texttt{run()} is called; orders match against the current
  book.
  \item Remaining player quotes rest; bots may trade against them.
  \item Unfilled resting quotes are cancelled; bot-to-bot trades may
  also occur.
  \item Next iteration begins with a fresh \texttt{TradingState}.
\end{enumerate}

Two implications.  First, you see the full book before acting, with no
information disadvantage relative to bots.  Second, resting quotes are
exposed to takers \emph{after} you submit, so passive MM accumulates
fills between iterations.

\subsection{Bot personalities}

Bots are not noise.  In P3, top teams identified distinct personalities:

\begin{itemize}
  \item \textbf{Market-maker bots} posting ``wall'' quotes at predictable
  distances from true price.  These walls are the most reliable
  fair-value signal: the ``wall mid'' (midpoint of innermost bid/ask
  walls) tracks the hidden true price
  closely~\citep{frankfurthedgehogs2025}.

  \item \textbf{Taker bots} crossing the spread.  Some are uninformed;
  others informed (e.g.\ ``Olivia'' in P3 bought the daily low and sold
  the daily high of Squid Ink).

  \item \textbf{Hidden taker bots} not visible in the book that fill
  resting orders priced attractively vs.\ a hidden fair value.  P3's
  Macarons had such a bot, filling $\approx$60\% of correctly-priced
  sells.
\end{itemize}

Identifying bot types is part of the pattern-recognition process
(\cref{ch:which_strategy_when}): once you know which bots are informed
vs noise, you can calibrate fair value and decide make vs take.

\begin{pitfallbox}[Do not assume bots are adversarial]
Prosperity bots are pre-programmed, not adaptive.  They do not learn
your strategy or front-run you, so exploitable patterns persist for the
round.  Flip side: \emph{taking} strategies have high backtester
fidelity (same bot behaviour replayed), but passive quoting is harder
to backtest because bots cannot react to your new quotes.
\end{pitfallbox}

%======================================================================
\section{A day-1 recipe}
\label{sec:day1_recipe}
%======================================================================

\begin{prosperitybox}
\textbf{Day-1 recipe: zero to profitable in 30 minutes.}

P3's Rainforest Resin had a fixed true value of 10{,}000.  Blank file
to profitable MM bot:

\begin{enumerate}
  \item \textbf{Install.} \texttt{pip install -U prosperity3bt}.
  \item \textbf{Write the stub.}  A \texttt{Trader} that buys at 9{,}999
  and sells at 10{,}001:
\begin{lstlisting}[numbers=none, frame=none,
  backgroundcolor=\color{orange!5}]
from datamodel import OrderDepth, TradingState, Order

class Trader:
    def run(self, state: TradingState):
        result = {}
        orders = []
        orders.append(Order("RAINFOREST_RESIN", 9999, 10))
        orders.append(Order("RAINFOREST_RESIN", 10001, -10))
        result["RAINFOREST_RESIN"] = orders
        return result, 0, ""
\end{lstlisting}
  \item \textbf{Run the backtester.}
  \texttt{prosperity3bt trader.py 0 -{}-vis}.
  \item \textbf{See PnL.}  The visualiser shows your trades and
  cumulative profit.  You are buying at 9{,}999 and selling at
  10{,}001, capturing a 2-unit spread each round trip.
  \item \textbf{Iterate.}  Widen quotes when inventory builds up
  (the Avellaneda--Stoikov skew from \cref{ch:market_making}).
  Tighten quotes when bots leave free money on the book (take
  any ask below 10{,}000 or any bid above 10{,}000).
\end{enumerate}

This 15-line bot will not win, but it establishes the full
pipeline (write, backtest, visualise, iterate) that top teams run
hundreds of times per round.  The Frankfurt Hedgehogs reported
Rainforest Resin alone contributed $\approx$39{,}000 SeaShells per round
from a simple variant.
\end{prosperitybox}

%======================================================================
\section{Pre-competition preparation checklist}
\label{sec:prosperity_prep}
%======================================================================

Top teams consistently attribute their placements to preparation done
\emph{before} the competition~\citep{frankfurthedgehogs2025}.  Before
Round~1 opens, have:

\begin{enumerate}
  \item \textbf{Backtester installed and tested.}  Run
  \texttt{prosperity3bt} on tutorial data; confirm log + visualiser.

  \item \textbf{Strategy skeletons for each archetype.}  Template code
  for: (a) flat-value MM, (b) trending MM (fair from wall mid), (c)
  basket arbitrage (synthetic price vs.\ residual), (d) options
  (Black--Scholes, IV extraction, delta hedge), (e) cross-exchange
  arbitrage (local vs.\ foreign after fees).  These need to
  \emph{exist}, not be tuned, so you can plug parameters in within the
  first hour.

  \item \textbf{Analysis notebooks.}  Load CSVs, plot prices, compute
  rolling mean/spread/autocorrelation, identify bot patterns.  Goal: full
  diagnostic in $<$15\,min per new product.

  \item \textbf{Logging framework.}  A parseable structured
  \texttt{print()} format: timestamp, product, position, fair estimate,
  orders, fills.  The Frankfurt Hedgehogs synced theirs with their
  dashboard's hover tooltip.

  \item \textbf{Study prior-year writeups.}  P2/P3 writeups (Stanford
  Cardinal, Linear Utility, Jasper Merle, Frankfurt Hedgehogs) expose
  what worked, what failed, and which bot patterns recur.
  Highest-ROI prep activity.

  \item \textbf{Read the entire wiki.}  Position limits, matching,
  libraries, the Round~2 \texttt{bid()} method; these are easy to miss
  if you only skim the API page.
\end{enumerate}

\begin{historybox}
IMC Prosperity began in 2022 as ``Prosperity~1,'' originally targeting
IMC's intern pipeline.  Its success prompted a global opening; by P3
(2025) it was among the largest student trading competitions, with
$>$12{,}000 teams from 100+ countries.  The design mirrors IMC's own
philosophy (market making, inventory management, robust
backtesting), so placing well is a credible signal to trading firms;
several top teams have received direct interview invitations from IMC
and competitors.
\end{historybox}

%======================================================================
\section{The Goldman Sachs parallel}
\label{sec:gs_parallel}
%======================================================================

\begin{gsbox}
Prosperity's \texttt{Trader} mirrors a production event-driven strategy
framework at a firm like Goldman Sachs:

\begin{itemize}
  \item \texttt{run()} is the \textbf{on-market-data callback}: when a
  tick arrives, the framework invokes your handler with current book
  state.
  \item \texttt{TradingState} is the \textbf{market-data snapshot}:
  positions, fills since last tick, order-book depth.
  \item The returned order list maps to a \textbf{smart order router}
  (SOR) or direct exchange gateway.
  \item A real desk runs microseconds, not 900\,ms, but the mental
  model is identical: receive state, compute fair, decide, submit.
  \item \texttt{traderData} has no production equivalent (real systems
  have databases and shared memory), but explicit state serialisation
  \emph{is} good practice: production strategies must handle restarts,
  and strats who design for explicit checkpointing save themselves
  hours after an unplanned restart.
\end{itemize}
\end{gsbox}

%======================================================================
\section{Key facts to memorise}
\label{sec:prosperity_keyfacts}
%======================================================================

\begin{enumerate}
  \item 5 rounds over $\approx$16 days, cumulative products: maintain
  prior-round strategies while adding new ones.

  \item Algorithm is a \texttt{Trader} class with
  \texttt{run(self, state)}: takes \texttt{TradingState}, returns
  \texttt{(orders\_dict, conversions, traderData)}.

  \item \texttt{OrderDepth.sell\_orders} values are \textbf{negative}.
  Negate for positive quantities.

  \item Position limits are all-or-nothing: aggregate could-exceed
  rejects \emph{every} order for that product.

  \item No persistent state between \texttt{run()} calls except
  \texttt{traderData} (max 50{,}000 chars).  Use JSON.

  \item \texttt{prosperity3bt} and the Jasper Merle visualiser are the
  community-standard tools.  Install before the competition.

  \item $\approx$6 recurring archetypes (flat-value MM, trending MM,
  signal-driven, basket, options, cross-exchange).  Identify first, then
  apply matching strategy.

  \item Zero latency: speed does not matter; \emph{fair-value} and
  \emph{inventory} determine PnL.

  \item Manual round is scored independently: free PnL, do not skip.

  \item Pre-competition: backtester, datamodel clone, strategy
  skeletons per archetype.
\end{enumerate}
